\subsection{Problem 1 and 2: Fixed Points, Periodic Points, and Newton's Method}

A discrete-time dynamical system describes the evolution of a state variable through successive iterations of a given map. In this  project's model, the state of the system \(n\) is described by the vector
\[
z_n = (x_n,y_n),
\]
where \(x_n\) and \(y_n\) denote the normalized densities of two interacting populations. The evolution of the system is determined by a map \(F\), such that
\[
z_{n+1}=F(z_n).
\]



\subsubsection{Fixed Points}

A fixed point of a map is a state that remains unchanged under one application of the map. Thus, a point \((x^*,y^*)\) is a fixed point of \(F\) if
\[
F(x^*,y^*)=(x^*,y^*).
\]


In this context, fixed points correspond to equilibrium states, where the reproductive dynamics and the interaction between the two populations are balanced. This reflects in a constant population density from one generation to the next. So, by definition, in order to get the fixed points, one has to find the roots of the system.

\subsubsection{Periodic Points}

Periodic points generalize the concept of fixed points. A point \((x,y)\) is said to be periodic with period \(p\) if it returns to itself after \(p\) applications of the map:
\[
F^p(x,y)=(x,y),
\]
For example,
\[
F^2(x,y)=F(F(x,y)),
\]
and
\[
F^3(x,y)=F(F(F(x,y))).
\]

A period-2 and 3 points therefore satisfy, respectively,
\[
F^2(x,y)=(x,y)
\]
and
\[
F^3(x,y)=(x,y),
\]
meaning that the orbit repeats after three iterations.

It is important, however, to distinguish between a point satisfying \(F^p(x,y)=(x,y)\) and a point having \emph{minimal period} \(p\). The minimal period is the smallest positive integer \(k\) such that
\[
F^k(x,y)=(x,y).
\]

For example, every fixed point satisfies
\[
F^2(x,y)=(x,y),
\]
but it should not be classified as a genuine period-2 point, because it already returns to itself after one iteration. Therefore, after solving the equation \(F^p(x,y)=(x,y)\), it is necessary to verify that the solution does not have a smaller period.

\subsection{Reformulation as a Root-Finding Problem}

The numerical computation of periodic points can be formulated as a nonlinear root-finding problem. This problem is solved with the aid of Newton's method and, for that, it is necessary to adapt the problem. For a given period \(p\), the function is defined as
\[
G_p(x,y)=F^p(x,y)-(x,y).
\]

To determine id a point is periodic, the following condition has to be met:
\[
G_p(x,y)=
\begin{pmatrix}
0 \\
0
\end{pmatrix}.
\]

This way, the problem is reduced strictly to finding the zeros of the function without the occurrence of repeated orbits. For \(p=2\), it solves

\[
G_2(x,y)=F^2(x,y)-(x,y)=0,
\]
while for \(p=3\), it solves
\[
G_3(x,y)=F^3(x,y)-(x,y)=0.
\]

This reformulation is useful and necessary because the iterated maps \(F^2\) and \(F^3\) generate nonlinear systems that are generally too complicated to solve analytically. A numerical method is therefore required.


